% ===== 正文部分 =====
\newgeometry{
	vmargin = 2cm,
	left = 3cm, 
	right = 2.5cm, 
	headheight	= 0.5cm,
	headsep		= 0.49cm,
	footskip	= 0.79cm, 
} % 仅本页生效的新边距，上下边距2cm，左3cm，右2.5cm

\chapter{\protect\hyperlink{:}{Introduction}}
\addtocontents{toc}{\protect\hypertarget{:}{}} %用\hypertarget实现反向跳转
\pagenumbering{arabic} % 切换阿拉伯数字
\setcounter{page}{1}   % 正文从1开始
\thispagestyle{fancy} % 强制当前页（第一页）使用 fancy 样式
% ===== 正文格式设置：Arial 11号,1.15倍行距，段后间距 10磅 =====
\setmainfont{Arial} % 设置正文字体为 Arial
\fontsize{11pt}{12.65pt}\selectfont % 11pt 字号，行距 11*1.15=12.65pt
\setlength{\parindent}{0pt} % 首行缩进（根据需要调整，通常英文为0）
\setlength{\parskip}{10pt} % 段后间距 10磅
\section{Background of Image Processing Technology}

Nowadays, with the continuous updating of the Internet and multimedia technologies, constantly upgrading of the electronic products such as mobile phones and video cameras, various images can be senn everywhere arounding us. Every day, we produce a large amount of image data, which is widely used in various fields such as scientific research, medical treatment, education and telecommunication. People's requirements for the quality of image data are getting higher and higher, so the process of image processing is indispensable, and the application field of image processing is becoming more and more extensive. Every contents of image processing are related to each other, common image processing techniques include image enhancement, image restoration, image recognition, image encoding, image compression, etc.
\subsection{Development of Image Processing Technology}

At th beginning, the aim of image processing is to improve image quality, input low-quality images and output improved high-quality images; it is now developing into new stages of automatic image recognition, such as face image recognition, satellite image positioning, and so on.

In the 1920s, digital image processing technology was first applied to the newspaper industry. Due to the need of information transmission in the newspaper industry, a submarine cable was connected from London, England to New York, USA, and the first photos was transmitted in the form of array. However, at that time, if digital image processing was not performed on the transmitted image, it will took 7 days for an image to be transmitted, but the time will reduce to 3 hours for transmission after digital image processing technology.

In the 1960s, a computer capable of performing image processing tasks was born. The emergence of the first image processing computer marked that the digital image processing technology beginning into a rapid development stage, using a computer can enable more advanced image processing. In these meantime, image processing technology has gradually been applied to medical imaging, astronomy and other fields. The most worth mentioning is the emergence of ct in the field of medical diagnosis, .

In the 1970s, with the development of computer technology and very large scale integrated circuits, image processing technology became more and more mature, and great achievements have been made in theoretical technology and practical applications.

In the 1980s, the field of application of digital image processing technology was further expanded, especially in the field of industrial inspection and geographic information. Digital image processing technology is applied to geographic information systems, in remote sensing technology, realized the image processing of the images transmitted from satellites.

In the 1990s, the birth of wavelet theory and wavelet transformation methods led to the rapid development of digital image processing technology, and do better in  the decomposition and reconstruction of digital images.

In the 21st century, the application range of digital image processing technology has been broadened again. It can be seen from the characteristics of current digital image processing technology that image processing technology has obvious advantages in processing large amount of information and high comprehensive images. Add to digital image processing technology also has high processing accuracy, flexibility, high reproducibility, wide range of applications and large information compression, the future application will be more extensive.
\subsection{Purpose of Image Processing}
In general, the purpose of processing images are three mainly as follows\cite{ref1}:

(1) Improve the real quality of the image, such as the brightness of the image, color conversion, enhancement/suppression of certain components of the image, geometric transformation of the image, etc., to improve the quality of the image.

(2) Extracting certain features or special information contained in the image, these extracted features or information often facilitate computer analysis of the image. The process of extracting features or information is preprocessing of pattern recognition or computer vision. The extracted features may include many aspects such as frequency domain features, grayscale or color features, boundary features, region features, texture features, shape features, topological features, and relationship structures.

(3) Transformation, encoding and compression of image data to facilitate image storage and transmission.
\subsection{Advantages of Digital Image Processing}
Digital images are processes that convert image signals into digital signals and process them using a computer. The main advantages are the following\cite{ref2}:

(1)High processing accuracy: In the image processing process, the simulated image can be converted into a two-dimensional array of any size for processing. 

(2)Flexibility: Any operations that can be expressed in mathematical formulas or logical relationships can be implemented with digital image processing.

(3)High reproducibility: A series of transformation operations of an image does not result in degradation of image quality.

(4)Wide range of applications: After the image is converted into a digitally encoded form, are always combined with the grayscale images represented by the two-dimensional array, so they can be processed by computer. 
\subsection{Application Fields of Digital Image Processing}
People obtain all kinds of information in life by image. Therefore, image processing must involve in all aspects of our lives. Here are a few examples of  the important applications field of digital image processing:  

(1)Aerospace technology: Image processing systems used in aircraft remote sensing and satellite remote sensing for image analysis.

(2)Biomedical engineering: In addition to the well-known CT technology, there are also microscopic image analysis.

(3)Telecommunication engineering field: The main development direction of current communication is multimedia communication combining sound, text, image and data. Specifically, the telephone, television, and computer are transmitted on a digital communication network in a triple play manner. Among them, image communication is the most complicated and difficult.

(4)Industrial engineering: automatic identification, inspection, sorting systems, etc. in automated assembly lines.

(5)Stylistic art field: The digital editing of TV pictures, the production of animation, electronic image games, textile handicraft design, fashion design and production, etc., gradually form a new art - computer art

(6)E-commerce field: Identity authentication, product anti-counterfeiting, watermark technology applied in E-commerce.
\subsection{Orthogonal Transformation in Digital Image Processing}
In the process of image processing, since the operation of the original image space may be complicated, in order to process and analyze the image efficiently and quickly, we can converted the original image space into another space in some form. Process the transformed image with the unique properties of another space, finally converts back to the original image space and achieves the desired image. This is the image transformation technology\cite{ref3}.

The orthogonal transformation of images is an important tool in image processing technology, and do various processing on the spectrum of an image. From a mathematical point of view, an image can be thought of as a two-dimensional signal, and the processed image is transformed into a two-dimensional array whose elements are called pixels and whose values are called grayscale. The orthogonal transform can make the image energy mainly concentrate on the low frequency components, and the edge and line information are concentrated on the high frequency components. Orthogonal transform processe the image by transforming the domain, which brings great convenience to the image processing work.
\section{Structure of Thesis}
The structure of this paper is mainly divided into three parts.

The first chapter introduces the background of image processing technology. The historical development process of image processing and the future development prospects are firstly studied. The technical advantages and broad application fields of image processing are analyzed. Based on these background, the orthogonal transformation in digital image processing is led out.

Tht second chapter introduces the four orthogonal transforms, Fourier transform, Discrete cosine transform, Walsh transform and Wavelet transform, which are widely used in digital image processing, learnd in principle and in nature. Nextly, the characteristics and practical application of these four kinds of transformations in digital image processing are analyzed.

The third chapter is based on the four orthogonal transformation learning in  theory, using Matlab simulation software to process the image. Four kinds of orthogonal transforms are programed in Matlab, and implement the image denoising and image compression functions in image processing by using different orthogonal transforms, and analyze and compare the results of image after processing.

%添加空白页，根据需要灵活取舍
% \clearpage
% \null
% \clearpage
